\documentclass[12pt,a4paper]{article}

% Greek language support
\usepackage[utf8]{inputenc}
\usepackage[LGR,T1]{fontenc}
\usepackage[greek,english]{babel}
\usepackage{alphabeta}

% Packages
\usepackage{geometry}
\usepackage{graphicx}
\usepackage{booktabs}
\usepackage{array}
\usepackage{float}
\usepackage{caption}
\usepackage{xcolor}
\usepackage{listings}
\usepackage{amsmath}
\usepackage{hyperref}
\usepackage{fancyhdr}
\usepackage{mdframed}

% Page settings
\geometry{left=2.5cm,right=2.5cm,top=2.5cm,bottom=2.5cm}

% Python code settings (simplified)
\definecolor{codegreen}{rgb}{0,0.6,0}
\definecolor{codegray}{rgb}{0.5,0.5,0.5}
\definecolor{backcolour}{rgb}{0.95,0.95,0.92}

\lstset{
    backgroundcolor=\color{backcolour},
    commentstyle=\color{codegreen},
    keywordstyle=\color{blue},
    basicstyle=\ttfamily\small,
    breaklines=true,
    numbers=left,
    numberstyle=\tiny\color{codegray},
    language=Python,
    frame=single
}

% Header/footer
\pagestyle{fancy}
\fancyhf{}
\rhead{AEM: 6931}
\lhead{Data Analysis 2025-26}
\cfoot{\thepage}

% Simple boxes
\mdfdefinestyle{infobox}{linecolor=blue,linewidth=1pt,backgroundcolor=blue!5}
\mdfdefinestyle{warnbox}{linecolor=orange,linewidth=1pt,backgroundcolor=orange!5}
\mdfdefinestyle{okbox}{linecolor=green,linewidth=1pt,backgroundcolor=green!5}

\title{
    \textbf{\textgreek{Ανάλυση Δεδομένων} 2025-26}\\[0.3cm]
    \large \textgreek{Ερώτημα 1: Δειγματοληψία και Καθαρισμός Δεδομένων}\\[0.2cm]
    \normalsize AEM: 6931
}
\author{}
\date{\textgreek{Μάρτιος} 2026}

\begin{document}
\maketitle
\tableofcontents
\newpage

\section{\textgreek{Εισαγωγή}}

\subsection{\textgreek{Σκοπός της Εργασίας}}

\textgreek{Η παρούσα εργασία αφορά την ανάλυση δεδομένων ποιότητας κρασιού. Οι στόχοι είναι:}

\begin{enumerate}
    \item \textgreek{Επιλογή τυχαίου δείγματος 80\% με seed τον ΑΕΜ}
    \item \textgreek{Έλεγχος ποιότητας δεδομένων}
    \item \textgreek{Καθαρισμός δεδομένων}
\end{enumerate}

\subsection{\textgreek{Τι είναι ο Καθαρισμός Δεδομένων;}}

\begin{mdframed}[style=infobox]
\textbf{\textgreek{Ορισμός:}} \textgreek{Ο καθαρισμός δεδομένων είναι η διαδικασία εντοπισμού και διόρθωσης λανθασμένων, ελλιπών, διπλότυπων ή ασυνεπών εγγραφών. Είναι κρίσιμο βήμα πριν από οποιαδήποτε ανάλυση.}
\end{mdframed}

\textgreek{Συνήθη προβλήματα:}
\begin{itemize}
    \item \textbf{\textgreek{Τυπογραφικά λάθη:}} \textgreek{π.χ. ``whte'' αντί ``white''}
    \item \textbf{\textgreek{Ελλείπουσες τιμές:}} \textgreek{Κενά πεδία}
    \item \textbf{\textgreek{Μη έγκυρες τιμές:}} \textgreek{π.χ. pH = 999}
    \item \textbf{\textgreek{Διπλοεγγραφές:}} \textgreek{Επαναλαμβανόμενες γραμμές}
\end{itemize}

\section{\textgreek{Περιγραφή Δεδομένων}}

\subsection{\textgreek{Μεταβλητές}}

\begin{table}[H]
\centering
\caption{\textgreek{Περιγραφή μεταβλητών}}
\begin{tabular}{clp{6cm}}
\toprule
\textbf{A/A} & \textbf{\textgreek{Μεταβλητή}} & \textbf{\textgreek{Περιγραφή}} \\
\midrule
1 & Fixed acidity & \textgreek{Σταθερή οξύτητα (g/L)} \\
2 & Volatile acidity & \textgreek{Πτητική οξύτητα (g/L)} \\
3 & Citric acid & \textgreek{Κιτρικό οξύ (g/L)} \\
4 & Residual sugar & \textgreek{Υπολειμματικά σάκχαρα (g/L)} \\
5 & Chlorides & \textgreek{Χλωριούχα άλατα (g/L)} \\
6 & Free SO2 & \textgreek{Ελεύθερο SO2 (mg/L)} \\
7 & Total SO2 & \textgreek{Συνολικό SO2 (mg/L)} \\
8 & Density & \textgreek{Πυκνότητα (g/cm3)} \\
9 & pH & \textgreek{Οξύτητα (2-14)} \\
10 & Sulphates & \textgreek{Θειικά άλατα (g/L)} \\
11 & Alcohol & \textgreek{Αλκοόλ (\%)} \\
12 & Quality & \textgreek{Βαθμολογία (0-10)} \\
13 & Wine\_type & \textgreek{red ή white} \\
\bottomrule
\end{tabular}
\end{table}

\textgreek{Αρχικό dataset:} \textbf{8.385} \textgreek{γραμμές} $\times$ \textbf{15} \textgreek{στήλες}

\section{\textgreek{Δειγματοληψία}}

\begin{mdframed}[style=warnbox]
\textbf{Seed:} \textgreek{Χρησιμοποιώντας τον ίδιο seed (ΑΕΜ=6931), παίρνουμε πάντα τα ίδια αποτελέσματα για αναπαραγωγιμότητα.}
\end{mdframed}

\begin{lstlisting}
import pandas as pd

df_original = pd.read_csv('Data Analysis_2026.csv')
AEM = 6931
df = df_original.sample(frac=0.80, random_state=AEM)
df = df.reset_index(drop=True)
\end{lstlisting}

\textgreek{Αποτέλεσμα:} 8.385 $\rightarrow$ \textbf{6.708} \textgreek{γραμμές (80\%)}

\section{\textgreek{Εντοπισμός Προβλημάτων}}

\subsection{\textgreek{Τυπογραφικά Λάθη (wine\_type)}}

\begin{table}[H]
\centering
\caption{\textgreek{Λάθη στη στήλη wine\_type}}
\begin{tabular}{lcc}
\toprule
\textbf{\textgreek{Λανθασμένη}} & \textbf{\textgreek{Πλήθος}} & \textbf{\textgreek{Σωστή}} \\
\midrule
Redd  & 75 & red \\
whit & 90 & white \\
r3d & 64 & red \\
WHITTE & 77 & white \\
whate & 39 & white \\
999, 9999, -999, ? & 215 & NaN \\
\midrule
\textbf{\textgreek{Σύνολο}} & \textbf{560} & \\
\bottomrule
\end{tabular}
\end{table}

\subsection{\textgreek{Μη Αριθμητικές Τιμές}}

\begin{table}[H]
\centering
\caption{\textgreek{Μη αριθμητικές τιμές ανά στήλη}}
\begin{tabular}{lc|lc}
\toprule
\textbf{\textgreek{Μεταβλητή}} & \textbf{N} & \textbf{\textgreek{Μεταβλητή}} & \textbf{N} \\
\midrule
fixed acidity & 87 & density & 53 \\
volatile acidity & 45 & pH & 151 \\
citric acid & 61 & sulphates & 63 \\
residual sugar & 56 & alcohol & 105 \\
chlorides & 57 & quality & 58 \\
free SO2 & 56 & & \\
total SO2 & 53 & \textbf{\textgreek{Σύνολο}} & \textbf{845} \\
\bottomrule
\end{tabular}
\end{table}

\subsection{\textgreek{Αδύνατες Τιμές}}

\textgreek{Τιμές 999, -999, 9999 εμφανίζονται ως κωδικοί ελλειπουσών τιμών.}

\begin{table}[H]
\centering
\caption{\textgreek{Λογικά όρια τιμών}}
\begin{tabular}{lcc|lcc}
\toprule
\textbf{\textgreek{Μεταβλητή}} & \textbf{Min} & \textbf{Max} & \textbf{\textgreek{Μεταβλητή}} & \textbf{Min} & \textbf{Max} \\
\midrule
Fixed acidity & 0 & 20 & Density & 0.9 & 1.1 \\
Volatile acidity & 0 & 5 & pH & 2 & 5 \\
Citric acid & 0 & 3 & Sulphates & 0 & 3 \\
Residual sugar & 0 & 100 & Alcohol & 5 & 20 \\
Chlorides & 0 & 1 & Quality & 0 & 10 \\
Free SO2 & 0 & 300 & & & \\
Total SO2 & 0 & 500 & & & \\
\bottomrule
\end{tabular}
\end{table}

\textgreek{Εκτός ορίων:} \textbf{2.037} \textgreek{τιμές}

\subsection{\textgreek{Διπλοεγγραφές}}

\begin{lstlisting}
duplicates = df.duplicated()
print(f"Duplicates: {duplicates.sum()}")
\end{lstlisting}

\begin{mdframed}[style=okbox]
\textgreek{Εντοπίστηκαν} \textbf{1.434} \textgreek{διπλοεγγραφές}
\end{mdframed}

\subsection{\textgreek{Ελλείπουσες Τιμές}}

\begin{table}[H]
\centering
\caption{Missing values}
\begin{tabular}{lcc|lcc}
\toprule
\textbf{\textgreek{Μεταβλητή}} & \textbf{N} & \textbf{\%} & \textbf{\textgreek{Μεταβλητή}} & \textbf{N} & \textbf{\%} \\
\midrule
fixed acidity & 574 & 10.9 & pH & 644 & 12.2 \\
volatile acidity & 542 & 10.3 & sulphates & 1 & 0.0 \\
density & 545 & 10.3 & alcohol & 600 & 11.4 \\
residual sugar & 134 & 2.5 & quality & 40 & 0.8 \\
\midrule
\multicolumn{3}{c}{} & \textbf{\textgreek{Σύνολο}} & \textbf{3.092} & \\
\bottomrule
\end{tabular}
\end{table}

\section{\textgreek{Αντιμετώπιση Προβλημάτων}}

\subsection{\textgreek{Στρατηγική}}

\begin{table}[H]
\centering
\caption{\textgreek{Στρατηγική αντιμετώπισης}}
\begin{tabular}{lp{7cm}}
\toprule
\textbf{\textgreek{Πρόβλημα}} & \textbf{\textgreek{Αντιμετώπιση}} \\
\midrule
\textgreek{Τυπογραφικά λάθη} & \textgreek{Διόρθωση ή NaN} \\
\textgreek{Μη αριθμητικές} & \textgreek{Μετατροπή σε NaN} \\
\textgreek{Αδύνατες τιμές} & \textgreek{Μετατροπή σε NaN} \\
\textgreek{Διπλοεγγραφές} & \textgreek{Αφαίρεση} \\
Missing values & \textgreek{Συμπλήρωση με διάμεσο ανά τύπο κρασιού} \\
\bottomrule
\end{tabular}
\end{table}

\subsection{\textgreek{Διόρθωση Τυπογραφικών}}

\begin{lstlisting}
corrections = {
    'Redd ': 'red', 'r3d': 'red',
    'WHITTE': 'white', 'whit': 'white',
    '999': None, '-999': None
}
for wrong, correct in corrections.items():
    df.loc[df['wine_type'] == wrong, 'wine_type'] = correct
\end{lstlisting}

\subsection{\textgreek{Γιατί Διάμεσος;}}

\begin{mdframed}[style=infobox]
\textgreek{Η διάμεσος είναι ανθεκτική σε ακραίες τιμές, σε αντίθεση με τον μέσο όρο. Υπολογίζεται ξεχωριστά για κόκκινα και λευκά κρασιά λόγω διαφορετικών χημικών χαρακτηριστικών.}
\end{mdframed}

\begin{table}[H]
\centering
\caption{\textgreek{Διάμεσοι ανά τύπο}}
\begin{tabular}{lcc}
\toprule
\textbf{\textgreek{Μεταβλητή}} & \textbf{\textgreek{Κόκκινο}} & \textbf{\textgreek{Λευκό}} \\
\midrule
Fixed acidity & 7.80 & 6.80 \\
Volatile acidity & 0.50 & 0.26 \\
Total SO2 & 40 & 133 \\
\bottomrule
\end{tabular}
\end{table}

\section{\textgreek{Αποτελέσματα}}

\subsection{\textgreek{Σύνοψη Προβλημάτων}}

\begin{table}[H]
\centering
\caption{\textgreek{Συνοπτικός πίνακας}}
\begin{tabular}{lcc}
\toprule
\textbf{\textgreek{Πρόβλημα}} & \textbf{\textgreek{Πλήθος}} & \textbf{\textgreek{Αντιμετώπιση}} \\
\midrule
\textgreek{Περιττές στήλες} & 2 & \textgreek{Αφαίρεση} \\
\textgreek{Τυπογραφικά wine\_type} & 560 & \textgreek{Διόρθωση} \\
\textgreek{Μη αριθμητικές} & 845 & NaN \\
\textgreek{Αδύνατες τιμές} & 2.037 & NaN \\
\textgreek{Διπλοεγγραφές} & 1.434 & \textgreek{Αφαίρεση} \\
Missing (\textgreek{αρχικά}) & 3.092 & \textgreek{Διάμεσος} \\
Missing (\textgreek{τελικά}) & 0 & - \\
\bottomrule
\end{tabular}
\end{table}

\subsection{\textgreek{Μέγεθος Dataset}}

\begin{table}[H]
\centering
\begin{tabular}{lcc}
\toprule
\textbf{\textgreek{Στάδιο}} & \textbf{\textgreek{Γραμμές}} & \textbf{\textgreek{Στήλες}} \\
\midrule
\textgreek{Αρχικό} & 8.385 & 15 \\
\textgreek{Μετά sample (80\%)} & 6.708 & 13 \\
\textgreek{Μετά duplicates} & 5.274 & 13 \\
\textbf{\textgreek{Τελικό}} & \textbf{5.274} & \textbf{13} \\
\bottomrule
\end{tabular}
\end{table}

\subsection{\textgreek{Κατανομή Τελικού}}

\begin{table}[H]
\centering
\begin{tabular}{lcc}
\toprule
\textbf{\textgreek{Τύπος}} & \textbf{\textgreek{Πλήθος}} & \textbf{\%} \\
\midrule
\textgreek{Λευκό} & 3.924 & 74.4 \\
\textgreek{Κόκκινο} & 1.350 & 25.6 \\
\midrule
\textbf{\textgreek{Σύνολο}} & \textbf{5.274} & \textbf{100} \\
\bottomrule
\end{tabular}
\end{table}

\section{\textgreek{Συμπεράσματα}}

\begin{enumerate}
    \item \textgreek{Τα πραγματικά δεδομένα περιέχουν πολλαπλά προβλήματα}
    \item \textgreek{Ο καθαρισμός είναι απαραίτητος πριν την ανάλυση}
    \item \textgreek{Κάθε πρόβλημα απαιτεί διαφορετική αντιμετώπιση}
    \item \textgreek{Οι αποφάσεις πρέπει να είναι τεκμηριωμένες}
\end{enumerate}

\textgreek{Το καθαρισμένο dataset είναι έτοιμο για περαιτέρω ανάλυση.}

\newpage
\appendix
\section{\textgreek{Πλήρης Κώδικας}}

\begin{lstlisting}
import pandas as pd
import numpy as np

# Load and sample
df = pd.read_csv('Data Analysis_2026.csv')
unnamed = [c for c in df.columns if 'Unnamed' in c]
df = df.drop(columns=unnamed)
df = df.sample(frac=0.80, random_state=6931).reset_index(drop=True)

# Fix wine_type typos
fixes = {'Redd ':'red', 'r3d':'red', 'WHITTE':'white', 
         'whit':'white', 'whate':'white',
         '999':np.nan, '9999':np.nan, '-999':np.nan, '?':np.nan}
for w, c in fixes.items():
    df.loc[df['wine_type']==w, 'wine_type'] = c

# Convert to numeric
num_cols = ['fixed acidity', 'volatile acidity', 'citric acid',
            'residual sugar', 'chlorides', 'free sulfur dioxide',
            'total sulfur dioxide', 'density', 'pH', 'sulphates',
            'alcohol', 'quality']
for col in num_cols:
    df[col] = pd.to_numeric(df[col], errors='coerce')

# Replace impossible values
limits = {'fixed acidity':(0,20), 'volatile acidity':(0,5),
          'citric acid':(0,3), 'residual sugar':(0,100),
          'chlorides':(0,1), 'free sulfur dioxide':(0,300),
          'total sulfur dioxide':(0,500), 'density':(0.9,1.1),
          'pH':(2,5), 'sulphates':(0,3), 'alcohol':(5,20),
          'quality':(0,10)}
for col, (mn, mx) in limits.items():
    df.loc[(df[col]<mn)|(df[col]>mx), col] = np.nan

# Remove duplicates
df = df.drop_duplicates().reset_index(drop=True)

# Fill missing values with median by wine type
df['wine_type'].fillna(df['wine_type'].mode()[0], inplace=True)
for col in num_cols:
    for wt in ['red','white']:
        mask = (df['wine_type']==wt) & df[col].isna()
        med = df.loc[df['wine_type']==wt, col].median()
        df.loc[mask, col] = med

# Save
df.to_csv('cleaned_wine_data_AEM_6931.csv', index=False)
\end{lstlisting}

\end{document}
